\section{Тема 15. Поверхность второго порядка в трёхмерном пространстве}

\subsection{Постановка задачи (вариант~4)}

Эллипсоид с полуосями $a,b,c$ задан каноническим уравнением
\[
  \frac{x^2}{a^2}+\frac{y^2}{b^2}+\frac{z^2}{c^2}=1.
\]
Определить, лежит ли заданная точка $(x,y,z)$ внутри, на
поверхности или вне эллипсоида.

\subsection{Математическое обоснование}

Введём функцию
\[
  f(x,y,z) = \frac{x^2}{a^2}+\frac{y^2}{b^2}+\frac{z^2}{c^2}.
\]
Тогда:
\begin{itemize}
  \item если $f=1$ (с учётом погрешности $\varepsilon$), точка
        лежит \emph{на поверхности} эллипсоида;
  \item если $f<1$, точка \emph{внутри};
  \item если $f>1$, точка \emph{вне} эллипсоида.
\end{itemize}

\subsection{Блок-схема алгоритма}

\begin{center}
\begin{tikzpicture}[
  node distance=10mm, start chain=going below,
  nodes={draw=black, top color=white, bottom color=lightgray!50,
         minimum width=5.6cm, align=center, font=\small},
  arrow/.style={->, >=Stealth, thick},
]
  \node[draw, rounded corners, on chain] {Start};
  \node[shape=trapezium, trapezium left angle=70,
        trapezium right angle=110, on chain, draw, join=by arrow]
       {$a,b,c$;\ $x,y,z$};
  \node[shape=rectangle, on chain, draw, join=by arrow]
       {$f := x^2/a^2{+}y^2/b^2{+}z^2/c^2$};
  \node[shape=diamond, aspect=2, on chain, draw, join=by arrow,
        minimum width=5.6cm]
       (c1) {$|f{-}1| < \varepsilon$\,?};
  \node[shape=trapezium, trapezium left angle=70,
        trapezium right angle=110, draw, right=2.2cm of c1]
       (s1) {«На поверхности»};
  \draw[arrow] (c1.east) -- node[above,font=\scriptsize]{$+$} (s1.west);
  \node[shape=diamond, aspect=2, on chain, draw,
        below=14mm of c1, minimum width=5.2cm]
       (c2) {$f < 1$\,?};
  \draw[arrow] (c1.south) -- (c2.north);
  \node[shape=trapezium, trapezium left angle=70,
        trapezium right angle=110, draw, right=2.2cm of c2]
       (s2) {«Внутри»};
  \draw[arrow] (c2.east) -- node[above,font=\scriptsize]{$+$} (s2.west);
  \node[shape=trapezium, trapezium left angle=70,
        trapezium right angle=110, on chain, draw, below=14mm of c2]
       (s3) {«Вне»};
  \draw[arrow] (c2.south) -- (s3.north);
  \node[draw, rounded corners, on chain, draw, below=10mm of s3]
       (end) {End};
  \draw[arrow] (s3.south) -- (end.north);
  \draw[arrow] (s1.south) |- (end.east);
  \draw[arrow] (s2.south) |- (end.east);
\end{tikzpicture}
\end{center}

\subsection{Текст программы}

\lstinputlisting[style=Cstyle]{../src/t15.c}

\subsection{Тестовые данные}

\begin{center}
\begin{tabular}{l|l}
\toprule
Эллипсоид, точка & Результат \\
\midrule
$a{=}b{=}c{=}1$, $(0,0,0)$ & Внутри ($f=0$) \\
$a{=}b{=}c{=}1$, $(1,0,0)$ & На поверхности ($f=1$) \\
$a{=}2,b{=}1,c{=}1$, $(3,0,0)$ & Вне ($f=2.25$) \\
\bottomrule
\end{tabular}
\end{center}

\textbf{Результат выполнения программы:}
\begin{lstlisting}[language={},basicstyle=\ttfamily\small,frame=single]
Полуоси эллипсоида (a b c): 1 1 1
Точка (x y z): 1 0 0
Точка лежит на поверхности эллипсоида (f=1.000000)
\end{lstlisting}
